
\documentclass{article}
\usepackage{colortbl}
\usepackage{makecell}
\usepackage{multirow}
\usepackage{supertabular}

\begin{document}

\newcounter{utterance}

\centering \large Interaction Transcript for game `cladder', experiment `full\_v1.5\_default', episode 2151 with Qwen3.6{-}27B.
\vspace{24pt}

{ \footnotesize  \setcounter{utterance}{1}
\setlength{\tabcolsep}{0pt}
\begin{supertabular}{c@{$\;$}|p{.15\linewidth}@{}p{.15\linewidth}p{.15\linewidth}p{.15\linewidth}p{.15\linewidth}p{.15\linewidth}}
    \# & \multicolumn{2}{c}{Player} && \multicolumn{2}{c}{Game Master} \\
    \hline

    \theutterance \stepcounter{utterance}  
    & & & \multicolumn{4}{p{0.6\linewidth}}{
        \cellcolor[rgb]{0.9,0.9,0.9}{
            \makecell[{{p{\linewidth}}}]{
                \texttt{\tiny{[P1$\langle$GM]}}
                \texttt{You are an expert in causal inference. The following question is not a typical commonsense query, but rather a meticulously designed question created by a professor specializing in causal inference, intended to assess the students' mastery of the course content.} \\
\texttt{Imagine a self{-}contained, hypothetical world with only the following conditions, and without any unmentioned factors or causal relationships: Gender has a direct effect on department competitiveness and admission status. Residency status has a direct effect on department competitiveness and admission status. Department competitiveness has a direct effect on admission status. Residency status is unobserved. The overall probability of male gender is 13\%. For individuals who are not male, the probability of admission acceptance is 61\%. For individuals who are male, the probability of admission acceptance is 65\%. Is admission acceptance less likely than admission rejection overall?} \\
\texttt{Start your answer with 'Yes' or 'No', followed by additional reasoning or evidence to support your explanation.} \\
            }
        }
    }
     \\ \\

    \theutterance \stepcounter{utterance}  
    & \multicolumn{4}{p{0.6\linewidth}}{
        \cellcolor[rgb]{0.561,0.737,0.561}{
            \makecell[{{p{\linewidth}}}]{
                \texttt{\tiny{[P1$\rangle$GM]}}
                \texttt{No} \\
\\ 
\texttt{**Reasoning:**} \\
\\ 
\texttt{To determine whether admission acceptance is less likely than admission rejection overall, we need to calculate the total probability of admission acceptance, denoted as \$P(A=1)\$, where \$A=1\$ represents acceptance and \$A=0\$ represents rejection.} \\
\\ 
\texttt{We are given the following information:} \\
\texttt{1.  The probability of being male is \$P(G=M) = 0.13\$.} \\
\texttt{2.  Consequently, the probability of not being male (female, in this binary context) is \$P(G=F) = 1 {-} 0.13 = 0.87\$.} \\
\texttt{3.  The probability of admission acceptance for individuals who are not male is \$P(A=1 | G=F) = 0.61\$.} \\
\texttt{4.  The probability of admission acceptance for individuals who are male is \$P(A=1 | G=M) = 0.65\$.} \\
\\ 
\texttt{Although the problem describes a complex causal structure involving unobserved residency status and department competitiveness, the question asks for the **overall** probability of admission acceptance. The law of total probability allows us to calculate the marginal probability of acceptance using only the observed gender distribution and the conditional probabilities of acceptance given gender. The unobserved variables and intermediate causal mechanisms do not change the marginal probabilities provided for the observed groups.} \\
\\ 
\texttt{Using the law of total probability:} \\
\texttt{\$\$ P(A=1) = P(A=1 | G=M) \textbackslash{}} \\
            }
        }
    }
    & & \\ \\

    \theutterance \stepcounter{utterance}  
    & & & \multicolumn{2}{p{0.3\linewidth}}{
        \cellcolor[rgb]{0.9,0.9,0.9}{
            \makecell[{{p{\linewidth}}}]{
                \texttt{\tiny{[GM$|$GM]}}
                \texttt{No} \\
            }
        }
    }
    & & \\ \\

    \theutterance \stepcounter{utterance}  
    & & & \multicolumn{2}{p{0.3\linewidth}}{
        \cellcolor[rgb]{0.9,0.9,0.9}{
            \makecell[{{p{\linewidth}}}]{
                \texttt{\tiny{[GM$|$GM]}}
                \texttt{no} \\
            }
        }
    }
    & & \\ \\

    \theutterance \stepcounter{utterance}  
    & & & \multicolumn{2}{p{0.3\linewidth}}{
        \cellcolor[rgb]{0.9,0.9,0.9}{
            \makecell[{{p{\linewidth}}}]{
                \texttt{\tiny{[GM$|$GM]}}
                \texttt{game\_result = WIN} \\
            }
        }
    }
    & & \\ \\

\end{supertabular}
}

\end{document}
